\section{Proofs}

\subsection{Auxiliary results}

\begin{definition}
 \label{def:filter}
 $\FF \subset \sigma(\Theta)$ is a 
 $\kappa$-complete filter over $\sigma(\Theta)$ if
 \begin{enumerate}
  \item $\Theta \in \FF$ and $\emptyset \notin \FF$.
	\item If $F \in \FF$ and $F \subset F' \in \sigma(\Theta)$, 
	      then $F' \in \FF$.
	\item If $\GG \subset \FF$, $|\GG| < \kappa$ and 
	      $\cap \GG \in \sigma(\Theta)$,
	      then $\cap \GG \in \FF$.
 \end{enumerate}
 In particular, if $\sigma(\Theta)=\PP(\Theta)$ and 
 $\kappa=\aleph_0$, then 
 $\FF$ is a filter over $\Theta$ \citep{Jech2013}.
\end{definition}

\begin{definition}
 $\AA(\phix) = \{H \in \sigma(\Theta): \phi(H|x) = 1\}$,
 $\RR(\phix) = \{H \in \sigma(\Theta): \phi(H|x) = 0\}$, and
 $\UU(\phix) = \left\{H \in \sigma(\Theta): \phi(H|x) = \half\right\}$.
\end{definition}

\begin{lemma}
 \label{thm:coherence-filter}
 A \summ, $\phi$, is 
 $\kappa$-coherent if and only if,
 for every $x \in \sX$,
 \begin{enumerate}
  \item $\AA(\phix)$ is a $\kappa$-filter 
  over $\sigma(\Theta)$, and
  \item $\RR(\phix) = \{H^{c}: H \in \AA(\phix)\}$.
 \end{enumerate}
\end{lemma}

\begin{proof}
 ($\Longrightarrow$) Let $\phi$ be 
 a $\kappa$-coherent \summ.
 First, we prove that $\AA(\phix)$ is 
 a $\kappa$-complete filter over $\sigma(\Theta)$.
 Since $\phi$ is $\kappa$-coherent, 
 \cref{def:coherence}.1 implies that
 $\Theta \in \AA(\phix)$.
 Also, it follows from 
 \cref{def:coherence}.(1 and 2) that
 $\emptyset \in \RR(\phix)$.
 Since, $\AA(\phix) \cap \RR(\phix) = \emptyset$,
 $\emptyset \notin \AA(\phix)$.
 Next, if $H \in \AA(\phix)$ and
 $H \subset H' \in \sigma(\Theta)$, then 
 \cref{def:coherence}.3 implies that
 $H' \in \AA(\phix)$.
 Finally, if $\GG \subset \AA(\phix)$, 
 $|\GG| < \kappa$ and
 $\cap \GG \in \sigma(\Theta)$, then
 it follows from \cref{def:coherence}.4 
 that $\cap \GG \in \AA(\phix)$.
 Hence, $\AA(\phix)$ is 
 a $\kappa$-complete filter over $\sigma(\Theta)$.
 It follows from \cref{def:coherence}.2 that
 $\RR(\phix) = \{H^{c}: H \in \AA(\phix)\}$.

 ($\Longleftarrow$) Let $\phi$ be a test such that, for every $x \in \sX$,
 $\AA(\phix)$ is a $\kappa$-complete filter over $\sigma(\Theta)$ and 
 also that $\RR(\phix) = \{H^{c}: H \in \AA(\phix)\}$.
 Take an arbitrary $x \in \sX$.
 Since $\AA(\phix)$ is a $\kappa$-complete filter, 
 $\Theta \in \AA(\phix)$.
 Since $\RR(\phix) = \{H^{c}: H \in \AA(\phix)\}$,
 conclude that $\emptyset = \Theta^c \in \RR(\phix)$.
 Therefore, $\phi$ satisfies \cref{def:coherence}.1.
 Also, since $\RR(\phix) = \{H^{c}: H \in \AA(\phix)\}$,
 $\phi$ satisfies \cref{def:coherence}.2.
 Next, let $H \subset H' \in \sigma(\Theta)$.
 To prove that $\phi$ satisfies \cref{def:coherence}.3,
 it is sufficient to show that
 $\phi(H|x) = 1 \rightarrow \phi(H'|x) = 1$ and
 $\phi(H'|x) = 0 \rightarrow \phi(H|x) = 0$.
 First, if $H \in \AA(\phix)$, then since 
 $\AA(\phix)$ is a $\kappa$-filter over $\sigma(\Theta)$,
 $H' \in \AA(\phix)$, that is, 
 $\phi(H|x) = 1 \rightarrow \phi(H'|x) = 1$.
 Second, if $H' \in \RR(\phix)$, then
 $H'^c \in \AA(\phix)$. Hence, since
 $\AA(\phix)$ is a $\kappa$-filter over $\sigma(\Theta)$,
 $H^c \in \AA(\phix)$ and $H \in \RR(\phix)$.
 That is, $\phi(H'|x) = 0 \rightarrow \phi(H|x) = 0$.
 From both cases, conclude that 
 $\phi$ satisfies \cref{def:coherence}.3.
 Finally, let $\GG \subset \AA(\phix)$ be such that 
 $|\GG| < \kappa$ and
 $\cap \GG \in \sigma(\Theta)$.
 Since $\AA(\phix)$ is a $\kappa$-complete filter,
 conclude that $\cap \GG \in \AA(\phix)$.
 Hence, $\phi$ satisfies \cref{def:coherence}.4.
 Since $\phi$ satisfies all of the conditions in
 \cref{def:coherence}, $\phi$ is $\kappa$-coherent.
\end{proof}

\begin{definition}
 \label{def:ultrafilter}
 $\FF \subset \sigma(\Theta)$ is a $\kappa$-complete ultrafilter over $\sigma(\Theta)$ if
 \begin{enumerate}
  \item $\FF$ is a a $\kappa$-complete filter over $\sigma(\Theta)$.
	\item For every $A \in \sigma(\Theta)$, 
	either $A \in \FF$ or $A^c \in \FF$.
 \end{enumerate}
 Note that, if $\sigma(\Theta)=\PP(\Theta)$,
 the power set of $\Theta$, and $\kappa=\aleph_0$,
 then $\FF$ is an ultrafilter over $\Theta$.
\end{definition}

\subsection{Proofs of \cref{cor:std_prob,thm:coherent_std_1}}

\begin{lemma}
 \label{cor:standard-ultrafilter}
  A \stdsumm, $\phi$, is
  $\kappa$-coherent if and only if,
  for every $x \in \sX$,
  $\AA(\phix)$ is
  a $\kappa$-complete ultrafilter over $\sigma(\Theta)$.
\end{lemma}

\begin{proof}
 Assume that $\phi$ is $\kappa$-coherent.
 It follows from \cref{thm:coherence-filter} that
 $\AA(\phix)$ is a $\kappa$-filter over $\sigma(\Theta)$ and
 $\RR(\phix) = \{H^c: H \in \AA(\phix)\}$.
 Since $\phi$ is a \stdsumm,
 $\AA(\phix)$ and $\RR(\phix)$ partition $\sigma(\Theta)$.
 Hence, $\AA(\phix)$ is a 
 $\kappa$-ultrafilter over $\sigma(\Theta)$.
 Next, assume that 
 $\AA(\phix)$ is a $\kappa$-ultrafilter over $\sigma(\Theta)$.
 In particular, $\AA(\phix)$ is a $\kappa$-filter over $\sigma(\Theta)$.
 Also, since $\phi$ is standard,
 $\RR(\phix) = \{H^c: H \in \AA(\phix)\}$.
 It follows from \cref{thm:coherence-filter} that
 $\phi$ is $\kappa$-coherent.
\end{proof}

\begin{proof}[Proof of \cref{cor:std_prob} and \cref{thm:coherent_std_1}]
 Follows directly from \cref{cor:standard-ultrafilter}.
\end{proof}

\subsection{Proof of \cref{{thm:coherence-nonstd}}}

\begin{lemma}
 \label{lemma:coherence_to_prob}
 If a \summ, $\phi$, is 
 $\aleph_0$-coherent, then
 for every $H_0 \in \sigma(\Theta)$ and 
 $x \in \sX$ there exists a probability, $\P_{H_{0}}$, such that:
 \begin{enumerate}
  \item $\P_{H_{0}}(H|x) = \phi(H|x)$, 
  for every $H$
  such that $\phi(H|x) \in \{0,1\}$,
  \item $\P_{H_{0}}(H_0|x) = 1$, 
  if $\phi(H_0|x) = \half$.
 \end{enumerate}
\end{lemma}

\begin{proof}  
 For each ultrafilter over $\sigma(\Theta)$, $\mathcal{F}$,
 let $\P_{\mathcal{F}}$ be a probability such that
 $\P_{\mathcal{F}}(H) = 1$ if 
 $H \in \mathcal{F}$ and
 $\P_{\mathcal{F}}(H) = 0$, otherwise.
 Assume $\phi$ is $\aleph_0$-coherent.
 It follows from \cref{thm:coherence-filter} that
 $\AA(\phix)$ is a filter over $\sigma(\Theta)$
 and that $\RR(\phix) = \{H: H^c \in \AA(\phix)\}$.
 Pick an arbitrary $H_0 \in \sigma(\Theta)$ and $x \in \sX$.
 If $\phix(H_0|x) \neq \half$, then let
 $\mathcal{F}$ be an arbitrary ultrafilter over $\sigma(\Theta)$
 such that $\AA(\phix) \subseteq \mathcal{F}$.
 The conditions of the lemma are satisfied choosing
 $\P_{H_0} := \P_{\mathcal{F}}$.
 If $\phi(H_0|x) = \half$, then let
 $\mathcal{F}$ be an
 ultrafilter over $\sigma(\Theta)$ such that
 $\AA(\phix) \subseteq \mathcal{F}$ and
 $H_0 \in \mathcal{F}$.
 The conditions of the lemma are satisfied by
 choosing $\P_{H_0} := \P_{\mathcal{F}}$.
\end{proof}

%\begin{theorem}
% \label{thm:agnostic_prob}
% A \summ, $\phi$, is 
% $\aleph_0$-coherent if and only if,
% for every $H_0 \in \sigma(\Theta)$ and 
% $x \in \sX$
% there exists a probability, $\P_{H_{0}}$, 
% such that:
% \begin{enumerate}
%  \item  $\P_{H_{0}}(H|x) = \phi(H|x)$ for every $H$
%  such that $\phi(H|x) \in \{0,1\}$,
%  \item $\P_{H_{0}}(H_0|x) \notin \{0,1\}$, if
%  $\phi(H_0|x) = \half$.
% \end{enumerate}
%\end{theorem}

%\begin{proof}[Proof of \cref{thm:agnostic_prob}] 
% ($\Longleftarrow$) Follows directly from %\cref{lemma:prob_to_coherence}.
% ($\Longrightarrow$) Let $\P_{H_0}$ and %$\P_{H_0^c}$ be
% such as in \cref{lemma:coherence_to_prob}. 
% The conditions of the theorem are satisfied choosing
% $\P^*_{H_0} = 0.5(\P_{H_0} + \P_{H_0^c})$.
%\end{proof} 

\begin{definition}
 Let $\sL$ be the linear space generated by
 $\{\I_H: H \in \sigma(\Theta)\}$ and
 $\sP$ be a family of probabilities over $\sigma(\Theta)$.
\end{definition}

\begin{lemma}
 \label{lemma:nonstd_1}
 For every $B \in \sigma(\Theta)$ and
 $X \in \sL$, $X \I_{B} \in \sL$
\end{lemma}

\begin{proof}
 There exist $(\alpha_i,A_i) \in \Re \times \sigma(\Theta)$,
 for $1 \leq i \leq n$, such that $X = \sum_{i} \alpha_i \I_{A_i}$.
 Note that
 \begin{align*}
  X \I_B 
  = \sum_{i} \alpha_i \I_{A_i} \I_B
  = \sum_{i} \alpha_i \I_{A_i \cap B} \in \sL
 \end{align*}
\end{proof}

\begin{definition}
 \label{def:pref-1}
 Let $X, Y \in \sL$.
 We say that $X \stackrel{\sP}{\preceq} Y$ if,
 for every $\P \in \sP$,
 $\E_{\P}[X] \leq \E_{\P}[Y]$, and
 $X \stackrel{\sP}{\prec} Y$ if
 $X \stackrel{\sP}{\preceq} Y$ and
 there exists $\P \in \sP$ such that
 $\E_{\P}[X] < \E_{\P}[Y]$.
 Also, $X \stackrel{*}{\prec} Y$ if
 $X \leq Y$ and there exists $\omega \in \Omega$
 such that $X(\omega) < Y(\omega)$.
 Note that $\stackrel{\sP}{\preceq}$ is 
 a weak partial order and
 both $\stackrel{\sP}{\prec}$ and
 $\stackrel{*}{\prec}$ are 
 strict partial orders.
\end{definition}

\begin{definition}
 \label{def:pref-2}
 We say that $X \preceq Y$ if either
 $X \leq Y$ or $X \stackrel{\sP}{\prec} Y$.
 Also, $X \ll Y$ if either
 $X \stackrel{*}{\prec} Y$ or $X \stackrel{\sP}{\prec} Y$.
\end{definition}

\begin{lemma}
 \label{lemma:prec_e}
 If $X \preceq Y$, then $X \stackrel{\sP}{\preceq} Y$.
\end{lemma}

\begin{lemma} 
 \label{lemma:nonstd_2}
 The following statements hold:
 \begin{enumerate}[leftmargin=*]
  \item For every $X \in \sL$, 
  $X \preceq X$. Also, there exist
  $X, Y \in \sL$ such that $X \prec Y$.
  
  \item If $Y-X \equiv Y'-X'$, then
  $X \preceq Y$ implies that $X' \preceq Y'$.
  
  \item If $X_j \preceq Y_j$ and 
  $\alpha_j \in \Re^{+}$ for every $j \in \{1,2\}$, then
  $\sum_{j=1}^2 \alpha_j X_j \preceq \sum_{j=1}^2 \alpha_j Y_j$.
  
  \item $X \prec Y$ if and only if $X \ll Y$.
  Also, if either $X \preceq Y$ and $Y \ll Z$ or
  $X \ll Y$ and $Y \preceq Z$, then $X \ll Z$.
  
  \item If $X \leq Y$, then $X \preceq Y$ and if
  $X \stackrel{*}{\prec} Y$, then $X \ll Y$.
 \end{enumerate}
\end{lemma}

\begin{proof}
 \begin{enumerate}[leftmargin=*]
  \item Since $X \leq X$, $X \preceq X$.
  Also, $\I_{\emptyset} \prec \I_{\Omega}$ since 
  $\I_{\emptyset} \stackrel{*}{\prec} \I_{\Omega}$ and
  neither $\I_{\Omega} \leq \I_{\emptyset}$ or
  $I_{\Omega} \stackrel{\sP}{\prec} \I_{\emptyset}$.
  
  \item $X \preceq Y$ implies that either 
  $X \leq Y$ or $X \stackrel{\sP}{\prec} Y$.
  If $X \leq Y$, since $Y-X \equiv Y'-X'$, 
  obtain $X' \leq Y'$ and, therefore, 
  $X' \preceq Y'$. Also,
  since $Y-X \equiv Y'-X'$, for every $\P \in \sP$,
  $\E_{\P}[X-Y] = \E_{\P}[X'-Y']$. That is,
  if $X \stackrel{\sP}{\prec} Y$, then 
  $X' \stackrel{\sP}{\prec} Y'$ and $X' \preceq Y'$.
  
  \item 
  Either $X_i \stackrel{\sP}{\prec} Y_i$ for some $i$ or
  $X_i \leq Y_i$ for every $i$.
  First, assume $X_i \stackrel{\sP}{\prec} Y_i$ for some $i$.
  Since $X_{3-i} \preceq Y_{3-i}$, it follows from
  \cref{lemma:prec_e} that 
  $X_{3-i} \stackrel{\sP}{\preceq} Y_{3-i}$.
  Since $X_i \stackrel{\sP}{\prec} Y_i$ and
  $X_{3-i} \stackrel{\sP}{\preceq} Y_{3-i}$, conclude that
  $\sum_{j=1}^2 \alpha_j X_j \stackrel{\sP}{\preceq}
  \sum_{j=1}^2 \alpha_j Y_j$, that is,
  $\sum_{j=1}^2 \alpha_j X_j \preceq
  \sum_{j=1}^2 \alpha_j Y_j$.
  Second, assume that $X_i \leq Y_i$ for every $i$. Conclude that
  $\sum_{j=1}^2 \alpha_j X_j \leq
  \sum_{j=1}^2 \alpha_j Y_j$, that is,
  $\sum_{j=1}^2 \alpha_j X_j \preceq
  \sum_{j=1}^2 \alpha_j Y_j$
  
  \item 
  First, we prove that $X \prec Y$ if and only if $X \ll Y$.
  If $X \preceq Y$, then either 
  $X \equiv Y$, or $X \stackrel{*}{\prec} Y$, 
  or $X \stackrel{\sP}{\prec} Y$. Since
  $\stackrel{*}{\prec}$ and $\stackrel{\sP}{\prec}$ are
  partial strict orders,
  $X \prec Y$ if and only if 
  $X \stackrel{*}{\prec} Y$ or
  $X \stackrel{\sP}{\prec} Y$, that is,
  $X \prec Y$ if and only if $X \ll Y$.
  
  Next, let either $X \ll Y$ and $Y \preceq Z$ or
  $X \preceq Y$ and $Y \ll Z$.
  There are three cases to consider.
  First, assume $X \stackrel{\sP}{\prec} Y$.
  It follows from \cref{lemma:prec_e} that
  $Y \stackrel{\sP}{\preceq} Z$ and, therefore,
  $X \stackrel{\sP}{\prec} Z$, that is,
  $X \ll Z$. Second, a symmetric chain of reasoning
  follows if  $Y \stackrel{\sP}{\prec} Z$.
  Finally, if neither $X \stackrel{\sP}{\preceq} Y$ or
  $Y \stackrel{\sP}{\preceq} Z$, then
  $X \leq Y$, $Y \leq Z$ and either 
  $X \stackrel{*}{\prec} Y$ or
  $Y \stackrel{*}{\prec} Z$. It follows that
  $X \stackrel{*}{\prec} Z$, that is,
  $X \ll Z$.
  
  \item Follows directly from definition.
 \end{enumerate}
\end{proof}

\begin{lemma}
 \label{lemma:coherence-nonstd_1}
 If $\phi$ is $\aleph_0$-coherent,
 then there exists a non-standard probability,
 $\P_x$, that represents $\phi$.
\end{lemma}

\begin{proof}
 Let $\phi$ be $\aleph_0$-coherent.
 It follows from \cref{lemma:coherence_to_prob} that
 for every $H_0 \in \sigma(\Theta)$ and $x \in \sX$
 there exists a probability, $\P_{H_{0}}$, such that
 $\P_{H_{0}}(H|x) = \phi(H|x)$ when $\phi(H|x) \in \{0,1\}$, and
 $\P_{H_{0}}(H_0|x) = 1$, if
 $\phi(H_0|x) = \half$.
 Let $\sP = \{\P_{H_0}: H_0 \in \sigma(\Theta)\}$,
 and $\preceq$ and $\ll$ be such as in
 \cref{def:pref-1,def:pref-2}. 
 Next, we show that $\ll$ satisfies
 four conditions that are 
 analogous to \cref{def:rep}.

 \begin{enumerate}
  \item Let $H_A^1$ and $H_A^2$ be such that
 $\phi(H_A^1|x) = \phi(H_A^1|x) = 1$.
 For every $\P \in \sP$,
 since $\P(H_A^1) = \P(H_A^2) = 1$,
 obtain that
 $\E_{\P}\left[\I_{H_A^1}\right] 
 = \E_{\P}\left[\I_{H_A^2}\right] = 1$.
 Hence, for every $\alpha \in (0,1)$,
 $\alpha\E_{P}\left[\alpha \I_{H_A^1}\right] 
 < \E_{P}\left[\I_{H_A^2}\right]$,
 that is,
 $\alpha\I_{H_A^1} 
 \stackrel{\sP}{\prec} \I_{H_A^2}$ and, therefore,
 $\alpha\I_{H_A^1} \ll \I_{H_A^2}$.
 
  \item Let $H_A$ and $H_U$ be such that
 $\phi(H_A|x) = 1$ and $\phi(H_U|x) = \half$.
 For every $\P \in \sP$, 
 $\P(H_A^c) = 0$ and
 $\E_{\P}[\I_{H_{A^c}}] = 0$.
 Furthermore, $\P_{H_U}(H_U) = 1$.
 Hence, for every $\alpha \in (0,1)$,
 $\E_{\P_{H_U}}\left[\alpha\I(H_U)\right]
 = \alpha > 0$.
 Conclude that $\I_{H_{A^c}} 
 \stackrel{\sP}{\prec} \alpha \I_{H_U}$ and,
 therefore, $\I_{H_{A^c}} 
 \ll \alpha \I_{H_U}$.
 Also, since $\phi$ is $\aleph_0$-coherent,
 it follows from invertibility that
 $\phi(H_U^c|x) = \half$. Hence, 
 from the previous considerations,
 $\I_{H_{A^c}} \ll \alpha \I_{H_U^c}$.

 \item Let $\phi(H_A|x) = 1$ and $\phi(H_R|x) = 0$.
 For every $\P \in \sP$,
 $\P(H_A) = 1$ and $\P(H_R) = 0$. Hence,
 for every $\alpha \in (0,1)$,
 $\E_\P[\I_{H_R}] = \E_\P[\I_{H_A^c}] = 0
 < \alpha = \E_\P[\alpha\I_{H_R^c}] 
 = \E_\P[\alpha\I_{H_A}]$. Hence,
 $\I_{H_R} \stackrel{\sP}{\prec} \alpha \I_{H_A}$
 and $\I_{H_A^c} \stackrel{\sP}{\prec} \alpha \I_{H_R^c}$ and, therefore,
 $\I_{H_R} \ll \alpha \I_{H_A}$
 and $\I_{H_A^c} \ll \alpha \I_{H_R^c}$.
 
 \item Let $\phi(H_R|x) = 0$ and
 $\phi(H_U|x) = \half$. For every $\P \in \sP$,
 $\P(H_R) = 0$ and $\E_\P[\I_{H_R}] = 0$. Also,
 $\P_{H_U}(H_U) = 1$. Hence,
 for every $\alpha \in (0,1)$,
 $\E_{\P_{H_U}}[\alpha \I_{H_U}] = \alpha > 0$.
 Conclude that $\I_{H_R} 
 \stackrel{\sP}{\prec} \alpha \I_{H_U}$ and, therefore,
 $\I_{H_R} \ll \alpha \I_{H_U}$.
 Also, since $\phi$ is $\aleph_0$-coherent,
 $\phi(H_U^c) = \half$ and, therefore,
 from the previous considerations,
 $\I_{H_R} \ll \alpha \I_{H_U^c}$.
 \end{enumerate}
 
 The proof follows from \cref{lemma:nonstd_1,lemma:nonstd_2} and
 \citet{Schervish2022}.
\end{proof}

\begin{lemma}
 \label{lemma:lemma:coherence-nonstd_2_1}
 If $\P_x$ represents a proper \summ, $\phi$,
 then for every $H \in \sigma(\Theta)$,
 $\phi(H|x) = 1 - \phi(H^c|x)$,
 that is, $\phi$ satisfies invertibility.
\end{lemma}

\begin{proof}
 There are three cases to consider:
 \begin{enumerate}
  \item $\left(\phi(H|x) = \half\right)$
  If $\phi(H^c|x) = 0$, then
  by taking $\alpha = 1$, obtain that
  $\P_x(H^c) < \P_x((H)^c)$, a contradiction.
  If $\phi(H^c|x) = 1$, then
  by taking $\alpha = 1$, obtain that
  $\P_x((H^c)^c) < \P_x(H)$, a contradiction.
  Conclude that $\phi(H^c|x) = \half$.

  \item $\left(\phi(H|x) = 0\right)$
  Since $\phi$ is proper, $\phi(\Theta|x) = 1$.
  If $\phi(H^c|x) = 0$, then by
  taking $\alpha = 0.1$ obtain that
  $\P_x(H) < 0.1 \P_x(\Omega) = 0.1$ and
  $\P_x(H^c) < 0.1 \P_x(\Omega) = 0.1$.
  Conclude that 
  $\P_x(H) + \P_x(H^c) = 0.2$, a contradiction.
  Therefore, $\phi(H^c|x) \neq 0$.
  Also, it follows from the previous item that
  $\phi(H^c|x) \neq \half$.
  Conclude that $\phi(H^c|x) = 1$.
 
  \item $\left(\phi(H|x) = 1\right)$.
  Since $\phi$ is proper, $\phi(\Theta|x) = 1$.
  If $\phi(H^c|x) = 1$, then by taking
  $\alpha = 0.9$ obtain that
  $\P_x(H|x) > 0.9\P_x(\Theta) = 0.9$ and
  $\P_x(H^c|x) > 0.9\P_x(\Theta) = 0.9$, that is,
  $\P_x(H|x) + \P_x(H^c|x) > 1.8$, a contradiction.
  Also, from the first item,
  $\phi(H^c|x) \neq \half$. Conclude that
  $\phi(H^c|x) = 0$.
 \end{enumerate}
\end{proof}

\begin{lemma}
 \label{lemma:lemma:coherence-nonstd_2_2}
 If $\P_x$ represents $\phi$ and
 $\P_x(H) \leq \P_x(H')$, then
 $\phi(H|x) \leq \phi(H'|x)$.
 Furthermore, if $\P_x$ represents $\phi$ and
 $H \subset H'$, then
 $\phi(H|x) \leq \phi(H'|x)$, that is,
 $\phi$ satisfies monotonicity.
\end{lemma}

\begin{proof}
 For the first part,
 assume $\P_x(H) \leq \P_x(H')$.
 There are three cases to consider:
 \begin{enumerate}
  \item $\phi(H|x) = 0$.  
  The result follows since
  $\phi(H'|x) \in \left\{0,\half,1\right\}$.

  \item $\phi(H|x) = \half$.
  Assume $\phi(H'|x) = 0$. Since
  $\P_x$ represents $\phi$,
  by taking $\alpha < 1$, obtain that
  $\P_x(H'|x) < \P_x(H)$, a contradiction.
  Hence, $\phi(H'|x) \geq \half = \phi(H|x)$.

  \item $\phi(H|x) = 1$.
  Assume $\phi(H'|x) \neq 1$.
  Since $\P_x$ represents $\phi$,
  by taking $\alpha < 1$, obtain that
  $1 - \P_x(H) < 1 - \P_x(H')$, that is,
  $\P_x(H) > \P_x(H')$, a contradiction.
  Hence $\phi(H'|x) = 1 \geq \phi(H|\x)$.
 \end{enumerate}
 For the second part, 
 let $H$ and $H'$ be such that
 $H \subset H'$. 
 Since $\P_x$ is a probability, 
 $\P_x(H) \leq \P_x(H')$. Hence,
 it follows from 
 the first part that
 $\phi(H|x) \leq \phi(H'|x)$.
\end{proof}

\begin{lemma}
 \label{lemma:lemma:coherence-nonstd_2_3}
 If $\P_x$ represents a \summ, $\phi$, and
 $\phi(H_i|x) = 1$, for every $1 \leq i \leq n$,
 then $\phi(\cap_{i=1}^n H_i|x) = 1$, that is,
 $\phi$ satisfies $\aleph_0$-consonance.
\end{lemma}

\begin{proof} 
 Let $H_* = \cap_{i=1}^n H_i$ and
 assume that $\phi(H_*|x) \neq 1$.
 Since $\P_x$ represents $\phi$, 
 it follows that
 \begin{align*}
  1-\P_x(H_i) &< n^{-1}(1 - \P_x(H_*)) 
  \nonumber \\
  \P_x(H_i^c) &< n^{-1} \P_x(H_*^c) 
  \nonumber \\
  \sum_{i=1}^n \P_x(H_i^c)
  &< \P_x(H_*^c) 
  \nonumber \\
  \P_x(\cup_{i=1}^n H_i^c)
  &< \P_x(H_*^c) 
  & \P_x(\cup_{i=1}^n H_i^c) 
  \leq \sum_{i=1}^n \P_x(H_i^c) 
  \nonumber \\
  \P_x(H_*)
  &> \P_x(H_*)
  & (\cup_{i=1}^n H_i^c)^c = H_*
 \end{align*}
 Conclude that $\phi(H_*|x) = 1$.
\end{proof}

\begin{lemma}
 \label{lemma:coherence-nonstd_2}
 If there exists a nonstandard probability,
 $\P_x$, that represents a
 proper \summ, $\phi$, then
 $\phi$ is $\aleph_0$-coherent.
\end{lemma}

\begin{proof}
 Follows directly from
 \cref{lemma:lemma:coherence-nonstd_2_1,lemma:lemma:coherence-nonstd_2_2,lemma:lemma:coherence-nonstd_2_3}.
\end{proof}

\begin{proof}[Proof of \cref{thm:coherence-nonstd}]
 Follows directly from
 \cref{lemma:coherence-nonstd_1,lemma:coherence-nonstd_2}
\end{proof}

\subsection{Proof of \cref{thm:second-countable}}

%\begin{lemma}
% \label{lemma:prob_1_filter}
% Let $\phi$ be \summ.
% Define
% \begin{align*}
%  \mathcal{F} = \{H \in \sigma(\Theta): 
%  \exists n \in \mathbb{N} \text{ and } 
%  H_1, H_2,\ldots, H_n \in \AA(\phix) 
%  \text{ s.t. } \cap_{i=1}^n H_i \subseteq H \}.
% \end{align*}
% If $\P$ is a probability such that
% $\P(H|x) = 1$ for every $H \in \AA(\phix)$, then
% $\P(H|x) = 1$ for every $H \in \mathcal{F}$.
%\end{lemma}

%\begin{proof}
% Let $H \in \mathcal{F}$.
% By construction, there exist 
% $H_1,\ldots,H_n \in \AA(\phix)$ such that
% $\cap_{i=1}^n H_i \subseteq H$.
% Let $\P$ be a probability such that
% $\P(H|x) = 1$ for every $H \in \AA(\phix)$.
% Since $\P$ is a probability,
% $\P(\cap_{i=1}^n H_i|x) = 1$. Hence,
% since $\cap_{i=1}^n H_i \subseteq H$, $\P(H|x) = 1$.
%\end{proof}

%\begin{lemma}
% \label{lemma:prob_to_coherence}
% If a \summ, $\phi$, is 
% not $\aleph_0$-coherent, then
% there exists no probability, $\P$, such that
% $\P(H|x) = \phi(H|x)$, for every $H$ such that
% $\phi(H|x) \in \{0,1\}$.
%\end{lemma}

%\begin{proof} 
% Assume $\phi$ is not
% $\aleph_0$-coherent. 
% It follows from \cref{thm:coherence-filter} that
% there exists $x \in \sX$ such that
% either $\AA(\phix)$ is not a filter 
% over $\sigma(\Theta)$ or
% $\RR(\phix) \neq \{H^c: H \in \AA(\phix)\}$.
 
% If $\AA(\phix)$ is not a filter, then
% let $\mathcal{F}$ be such as in %\cref{lemma:prob_1_filter}.
% Since $\mathcal{F}$ is a filter,
% there exists $H_0 \in \mathcal{F}$ such that
% $H_0 \notin \AA(\phix)$. Using 
% \cref{lemma:prob_1_filter} obtain that there %exists 
% no probability, $\P$, such that $\P(H|x) = 1$
% for every $H \in \AA(\phi_x)$ and 
% $\P(H_0|x) \neq 1$.

% If $\RR(\phix) \neq \{H^c: H \in \AA(\phix)\}$, then
% there exists $H_0 \in \sigma(\Theta)$ such that
% $\phi(H_0|x) \in \{0,1\}$ and $\phi(H_0^c|x) \neq 1-\phi(H_0|x)$.
% If $\phi(H_0|x) = \phi(H_0^c|x)$, then there exists 
% no probability, $\P$, such that
% $\P(H_0|x) = 1-\phi(H|x) = \P(H_0^c|x)$.
% If $\phi(H_0^c|x) = \half$, then for every probability, $\P$,
% such that $\P(H_0|x) = 1-\phi(H_0|x)$, obtain that
% $\P(H_0^c|x) = 1-\P(H_0|x) \in \{0,1\}$.
%\end{proof}

\subsection{Proof of \cref{thm:min_u}}

\begin{proof}[Proof of \cref{thm:min_u}]
 Let 
 \begin{align*}
  U_1(\phi, \theta) &= \I(\theta \in P_{\phi}(X)) - \lambda \mu(P_{\phi}(X)), \text{ and } \\
  U_2(\phi, \theta) &= \I(\theta \in N_{\phi}(X)) - \lambda \mu(N_{\phi}(X)).
 \end{align*}
 It follows from \citet{Schervish2012} that
 $\phi$ maximizes $U_1$ and $U_2$ 
 if and only if $P_{\phi}(X)$ and $N_{\phi}(X)$
 a.s. satisfy
 \begin{align}
  \label{eq:min_u_1}
  P_{\phi}(X) = N_{\phi}(X) =
  R_x :=
  \left\{\theta \in \Theta: 
  \frac{d\P(\theta|X)}{d\mu} > (1+\lambda)^{-1} \right\}.
 \end{align}
 Since $U = U_1 + U_2 -\odot \cdot \rho(U_{\phi})$
 and $\odot \cdot \rho(U_{\phi})$ is infinitesimal,
 conlude that \cref{eq:min_u_1} is 
 a necessary condition for
 $\phi$ to maximize $U$.
 Furthermore, it follows from \cref{def:region_p_n} that \cref{eq:min_u_1} holds 
 if and only if
 $\phi(H|x) \geq \half$ when
 $H \cap R_x \neq \emptyset$ and
 $\phi(H|x) \leq \half$ when
 $H^c \cap R_x \neq \emptyset$.
 From these conditions, one obtains that
 $\phi(H|x) = \half$ whenever
 $H \cap R_x \neq \emptyset$ and
 $H^c \cap R_x \neq \emptyset$.
 Since the \summ\ based on $R_x$, $\phi^*$,
 obtains $\phi(H|x) = \half$ if and only if
 $H \cap R_x \neq \emptyset$ and
 $H^c \cap R_x \neq \emptyset$, conclude that
 $\phi^*$ minimizes $\rho(A_{\phi}(X))$ among
 the $\phi$ that satisfy \cref{eq:min_u_1}.
 That is, $\phi^*$ maximizes $U$.
\end{proof}